% =====================================================================
%  FICHE 11 — THEME 5 — Corrigé FACILE — Piles, Nernst, Faraday
% =====================================================================
\documentclass[11pt,a4paper]{article}
\usepackage[utf8]{inputenc}
\usepackage[T1]{fontenc}
\usepackage{lmodern}
\usepackage[french]{babel}
\usepackage{amsmath,amssymb}
\usepackage[version=4]{mhchem}
\usepackage{siunitx}
\sisetup{output-decimal-marker={,},per-mode=symbol}
\usepackage{xcolor}
\usepackage{array}
\usepackage{booktabs}
\usepackage{enumitem}
\usepackage[most]{tcolorbox}
\usepackage[a4paper,margin=2.1cm,headheight=26pt,headsep=10pt]{geometry}
\usepackage{fancyhdr}
\usepackage[hidelinks]{hyperref}

\definecolor{primary}{RGB}{146,95,160}
\definecolor{primarydark}{RGB}{92,52,104}
\definecolor{corrcol}{RGB}{0,135,90}
\newcommand{\NO}{\ensuremath{\mathrm{N.O.}}}
\newcommand{\Ered}{\ensuremath{E_{\mathrm{r}}^{\circ}}}
\newcommand{\fem}{\ensuremath{\Delta E}}

\newcounter{exo}
\newtcolorbox{corr}[1][]{enhanced,breakable,boxrule=0.9pt,arc=2pt,colback=corrcol!4,
  colframe=corrcol,fonttitle=\bfseries,coltitle=white,left=3mm,right=3mm,top=2.2mm,bottom=2.2mm,
  title={Corrigé~\refstepcounter{exo}\theexo\ifx\relax#1\relax\relax\else\ \textmd{--- \itshape #1}\fi}}

\pagestyle{fancy}\fancyhf{}
\fancyhead[L]{\small\textcolor{primarydark}{\textbf{Fiche 11 · Thème 5} — Corrigé}}
\fancyhead[R]{\small\textcolor{primarydark}{\textsc{Facile}}}
\fancyfoot[C]{\small\thepage}
\renewcommand{\headrulewidth}{0.8pt}
\renewcommand{\headrule}{\hbox to\headwidth{\color{primary}\leaders\hrule height \headrulewidth\hfill}}

\begin{document}
\begin{center}
{\LARGE\bfseries\color{primarydark}Thème 5 — Corrigé Facile}\\[-2pt]
{\color{corrcol}\rule{\textwidth}{1pt}}
\end{center}

\begin{corr}[anode ou cathode ?]
Le couple de \Ered\ le plus faible est l'anode, le plus élevé la cathode.
\begin{enumerate}[label=(\alph*),leftmargin=2em,topsep=1pt,itemsep=1.5pt]
  \item $\Ered(\ce{Zn^{2+}}/\ce{Zn})=-0{,}76 < \Ered(\ce{Cu^{2+}}/\ce{Cu})=+0{,}34$ : le \textbf{zinc est
        l'anode} ($-$), le \textbf{cuivre la cathode} ($+$).
  \item Anode (oxydation) : $\ce{Zn -> Zn^{2+} + 2 e^-}$. Cathode (réduction) : $\ce{Cu^{2+} + 2 e^- -> Cu}$.
\end{enumerate}
\end{corr}

\begin{corr}[circulation des porteurs de charge]
\begin{enumerate}[label=(\alph*),leftmargin=2em,topsep=1pt,itemsep=1.5pt]
  \item Les électrons vont de l'\textbf{anode ($-$) vers la cathode ($+$)} dans le fil (l'anode les
        produit, la cathode les consomme). Le courant conventionnel $i$ va en sens inverse.
  \item Le pont salin \textbf{ferme le circuit} et maintient l'\textbf{électroneutralité} de chaque
        compartiment. Les \textbf{anions} migrent vers l'\textbf{anode} (et les cations vers la cathode).
\end{enumerate}
\end{corr}

\begin{corr}[écrire la notation conventionnelle]
Anode (Zn) à gauche, cathode (Ag) à droite ; \textbar\ pour l'interface, $\|$ pour le pont salin :
\[
\boxed{\ce{(-) Zn | Zn^{2+} || Ag+ | Ag (+)}}
\]
\end{corr}

\begin{corr}[force électromotrice standard]
\[
\fem^{\circ} = \Ered(\text{cathode}) - \Ered(\text{anode}) = \Ered(\ce{Ag+}/\ce{Ag}) - \Ered(\ce{Zn^{2+}}/\ce{Zn})
= 0{,}80 - (-0{,}76) = \boxed{+1{,}56\ \mathrm{V}}.
\]
$\fem^{\circ}>0$ : la pile débite bien spontanément.
\end{corr}

\begin{corr}[première approche de la loi de Faraday]
\begin{enumerate}[label=(\alph*),leftmargin=2em,topsep=1pt,itemsep=1.5pt]
  \item $Q = I\,t = 2{,}0 \times 1930 = \SI{3860}{C}$.
  \item $n_{\ce{e^-}} = \dfrac{Q}{F} = \dfrac{3860}{96485} \simeq \SI{0.0400}{mol}$ d'électrons.
\end{enumerate}
C'est l'étape clé : $Q/F$ donne les moles d'électrons, dont on déduira ensuite la masse déposée.
\end{corr}

\end{document}
